\documentclass[11pt]{article}
\usepackage[margin=1in]{geometry}
\usepackage{amsmath,amssymb,amsthm,mathtools}
\usepackage[T1]{fontenc}
\usepackage[expansion=false]{microtype}
\usepackage[hidelinks]{hyperref}
\newtheorem{theorem}{Theorem}[section]
\newtheorem{lemma}[theorem]{Lemma}
\newtheorem{proposition}[theorem]{Proposition}
\newtheorem{remark}[theorem]{Remark}
\newcommand{\ddc}{\mathrm{d}\mathrm{d}^c}
\newcommand{\dc}{\mathrm{d}^c}
\newcommand{\MA}{\operatorname{MA}}
\newcommand{\PSH}{\operatorname{PSH}}
\newcommand{\supp}{\operatorname{supp}}
\newcommand{\eps}{\varepsilon}
\newcommand{\Om}{\Omega}
\title{Nonclosed Monge--Amp\`ere equations on cuspidal products}
\author{}
\date{}
\begin{document}
\maketitle
\begin{abstract}
In every complex dimension $n\ge2$, we construct a compact nonnormal
unibranch space with a genuinely nonclosed Hermitian background in the local
embedding sense. Its normalization is a flat
complex torus, and the pullback background degenerates along the cusp divisor.
An explicit bounded reference repairs it to a positive metric whose powers
are all $\ddc$-closed, because its torsion is supported on two factors.
Fourier inversion and a contraction solve the nonlinear equation for an open
family of positive Sobolev densities depending on all complex variables.
A decreasing-convolution energy argument gives uniqueness against every
bounded canonical competitor. We identify the normalization constant, prove
continuity of the entire potential across the cusp with exponent $1/3$, and
verify independence under every resolution refinement. The result concerns
this product family and bounded ordinary type; it does not classify arbitrary
unibranch singularities or potentials with poles.
\end{abstract}

\section{Conventions and result}

Throughout,
\[
 \ddc=\frac{i}{2\pi}\partial\bar\partial,
 \qquad \dc=\frac{i}{4\pi}(\bar\partial-\partial).
\]
Volume forms are $n$th powers without a factor $1/n!$.
All potentials have their genuine canonical upper values. On a complex space
with a smooth local embedding background $h$, its disk definition requires
$u\circ f+p_f$ to be subharmonic or identically $-\infty$ for every
holomorphic disk $f$, where $\ddc p_f=f^*h$ locally. Ordinary type zero is the
entire class of bounded genuine $h$-plurisubharmonic functions.

The extended measure is computed on a smooth model and pushed down.
For bounded potentials its local products are Bedford--Taylor products with
the smooth background corrections. We prove the required independence of
the smooth model in Section~\ref{sec:refinements}.

Let $E_1,\ldots,E_n$ be copies of the square elliptic curve
$E=\mathbb C/(\mathbb Z+i\mathbb Z)$, let $p\in E_1$, and let $C$ be the cuspidal
pinching at $p$ constructed in Section~\ref{sec:geometry}. Set
\[
 X=C\times\prod_{j=2}^n E_j,\qquad Y=\prod_{j=1}^n E_j,\qquad
 \pi=\pi_1\times\mathrm{id}:Y\longrightarrow X.
\]
The smooth forms and constants in the following statement are explicitly
constructed before any equation is solved.

\begin{theorem}\label{thm:main}
For every $n\ge2$ there are a smooth local embedding Hermitian background $h$ on $X$ with
$\mathrm{d}h\ne0$, a bounded continuous genuine $h$-psh reference descending
from a smooth function $v_0$ on $Y$, and a flat K\"ahler form $K_0$ on $Y$,
such that
\[
 H=\pi^*h,\qquad k=H+\ddc v_0\ge\tfrac12K_0,
 \qquad \ddc(k^r)=\ddc(H^r)=0\quad(r\ge0).
\]
The form $H$ has rank $n-1$ along $D=\{p\}\times\prod_{j=2}^nE_j$.
For each integer $s\ge n+3$, the construction allows a sufficiently small
nonzero background perturbation and an explicit $\delta>0$ with the
following property. Put
\[
 \Om=K_0^n,\qquad M=\int_Y k^n,\qquad f_\eps=\frac{k^n}{\Om}.
\]
Let $f>0$ be real and in $H^s(Y)$, and assume
\begin{equation}\label{eq:data-condition}
 \left\|\frac{M}{\int_Yf\Om}f-f_\eps\right\|_{H^s}\le\delta.
\end{equation}
Then for $\nu=\pi_*(f\Om)$ there is a unique supremum-normalized solution
in ordinary type zero of
\[
 \MA^{\mathrm{ext}}_h(u)=c\nu,\qquad
 c=\frac{M}{\nu(X)}>0,\qquad \sup_Xu=0.
\]
Its full potential is locally H\"older continuous with exponent $1/3$ in
local ambient coordinates, including along the cusp. Its pullback is
$v_0+w$ up to a constant with $w\in H^{s+2}(Y)$; in particular the solution
is $C^2$ on the regular locus. Every bounded genuine potential in the
specified type has finite total Monge--Amp\`ere mass $M$.
Both the type and measure are independent of every resolution and common
refinement.
\end{theorem}

The data condition is invariant under positive rescaling of $f$. After fixing
its integral, strict inequality defines an open set in the full Sobolev
space. Densities need not separate into functions of the factors.

\section{The cusp and a nonclosed background}\label{sec:geometry}

Choose a coordinate $z$ on a small disk centered at $p$ in $E_1$. Replace
the disk by the cusp $y^2=x^3$, with normalization
\[
 z\longmapsto (x,y)=(z^2,z^3),
\]
and glue its puncture to $E_1\setminus\{p\}$ by this map.
The resulting compact curve $C$ has the same underlying topological space
as $E_1$ and one singular point $q$. Its local ring is
$\mathbb C\{z^2,z^3\}$, a domain with one analytic branch. The normalization
$\pi_1:E_1\to C$ is finite: locally $z$ is integral, with $z^2=x$, and
the map is proper. It is bijective of degree one, hence a homeomorphism.
Consequently $X$ is a compact, connected, reduced, locally irreducible
complex space of pure dimension $n$, with singular locus
$\{q\}\times\prod_{j=2}^nE_j$.

Choose the flat forms $\kappa_j=(i/2\pi)\,\mathrm{d}z_j\wedge
\mathrm{d}\bar z_j$ on $E_j$ and a local embedding
Hermitian form $h_C$ on $C$ which near $q$ equals
\[
 \ddc(|x|^2+|y|^2)|_C.
\]
Away from $q$ one may glue positive curve metrics by smooth partitions;
the gluing occurs on the regular curve and admits positive ambient
extensions. Near the singular point the displayed extension is already
strictly positive. The smooth pullback $\eta=\pi_1^*h_C$ satisfies
\[
 \eta=\frac{i}{2\pi}(4|z|^2+9|z|^4)\,\mathrm{d}z\wedge
                  \mathrm{d}\bar z
 \quad\hbox{near }p.
\]
It is nonnegative, positive away from $p$, and vanishes at $p$.
Write $\eta=\rho\kappa_1$ and define
\[
 \lambda=\frac{\int_{E_1}\eta}{\int_{E_1}\kappa_1}>0,
 \qquad h_0=h_C+\sum_{j=2}^n\kappa_j,
\]
\[
 H_0=\pi^*h_0=\eta+\sum_{j=2}^n\kappa_j,
 \qquad K_0=\lambda\kappa_1+\sum_{j=2}^n\kappa_j.
\]
Product pullbacks are suppressed in this notation. The form $H_0$ is
closed and has rank $n-1$ on $D$.

Take a smooth real function $A$ on $C$ supported away from $q$ and a
smooth real function $B$ on $E_2$. The real $(1,1)$-form
\begin{equation}\label{eq:alpha}
 \alpha=\partial A\wedge\bar\partial B+
                  \bar\partial A\wedge\partial B
\end{equation}
is smooth in the local embedding sense and vanishes near the singular locus.
It is pulled back from the first two factors only.
Indeed,
\[
 \alpha=\partial(A\bar\partial B)+\bar\partial(A\partial B),
 \qquad \ddc\alpha=0,
\]
where the two terms involving $A\partial\bar\partial B$ cancel.
On the other hand,
\[
 \partial\alpha=-\partial A\wedge\partial\bar\partial B
                      +\partial\bar\partial A\wedge\partial B.
\]
Choose a regular point where $\partial A=0$ and
$\partial\bar\partial A\ne0$, and a point of $E_2$ where
$\partial B\ne0$. Local quadratic and linear functions with cutoffs
provide such choices. Thus $\mathrm{d}\alpha\ne0$.

Denote its pullback to $Y$ by $\widetilde\alpha$. Both constants
\[
 r_H=\sup_{Y\setminus D}|\widetilde\alpha|_{H_0,\mathrm{op}},
 \qquad r_K=\sup_Y|\widetilde\alpha|_{K_0,\mathrm{op}}
\]
are finite; the first uses its vanishing near $D$.
Choose $\eps\ne0$ with $|\eps|r_H<1/2$ and $|\eps|r_K<1/2$.
Then
\[
 h=h_0+\eps\alpha,\qquad H=H_0+\eps\widetilde\alpha
\]
defines a positive local embedding background with
$\mathrm{d}h=\eps\mathrm{d}\alpha\ne0$. It equals $h_0$ near the cusp.
Upstairs $H$ remains degenerate on $D$, and $\ddc H=0$.

The smooth function $\lambda-\rho$ has zero $\kappa_1$-integral.
Fourier division on $E_1$ gives a unique smooth mean-zero real $v$ with
\[
 \ddc v=(\lambda-\rho)\kappa_1.
\]
Each nonzero Fourier coefficient is divided by a nonzero negative
Laplacian eigenvalue comparable to the squared frequency; rapid decay
preserves smoothness. With $v_0=\mathrm{pr}_1^*v$,
\begin{equation}\label{eq:repaired}
 k=H+\ddc v_0=K_0+\eps\widetilde\alpha\ge\tfrac12K_0,
 \qquad \ddc k=0.
\end{equation}
This proves the positive repair rather than assuming positivity of a
normalization pullback.

Moreover $\mathrm{d}k$ and $\dc k$ are pulled back from the smooth complex
surface $E_1\times E_2$. They are three-forms, so their wedge is a
six-form on a real four-dimensional source and vanishes identically.
The same holds for $H$, since $H_0$ is closed. Thus
\[
 \ddc k=\ddc H=0,\qquad
 \mathrm{d}k\wedge\dc k=\mathrm{d}H\wedge\dc H=0.
\]
The Leibniz rule gives $\ddc(k^r)=\ddc(H^r)=0$ for every $r\ge0$.
This additional torsion identity is essential in dimension greater than two;
a generic pluriclosed metric need not have $\ddc$-closed powers.

\section{Fourier inversion and nonlinear existence}\label{sec:existence}

Use Fourier Sobolev norms on the real $2n$-dimensional torus $Y$, and let
$H^r_0$ consist of real functions of zero $\Om$-integral. In this section
write $\alpha$ for $\widetilde\alpha$ and fix an integer $s\ge n+3$.
Define
\[
 L_0w=\frac{nK_0^{n-1}\wedge\ddc w}{\Om},\qquad
 L_\eps w=\frac{nk^{n-1}\wedge\ddc w}{\Om}.
\]
Write $b_1=\lambda$ and $b_j=1$ for $j\ge2$. The operator $L_0$ is the
constant coefficient Laplacian
\[
 L_0=\sum_{j=1}^n b_j^{-1}\partial_{z_j}\partial_{\bar z_j},
 \qquad \ell(m)=-\pi^2\sum_{j=1}^n b_j^{-1}
                       (m_{2j-1}^2+m_{2j}^2).
\]
Its nonzero Fourier eigenvalues are negative and comparable in magnitude to $|m|^2$.
Fourier division therefore gives a bounded inverse
$L_0^{-1}:H^s_0\to H^{s+2}_0$. Put
\[
 A_0=\sup_{m\ne0}\frac{1+|m|^2}{|\ell(m)|}
                  \le\frac{2\max_j b_j}{\pi^2}.
\]
The perturbed operator preserves zero mean, since
\[
 \int_Y L_\eps w\Om=n\int_Yk^{n-1}\wedge\ddc w
                         =n\int_Yw\ddc(k^{n-1})=0.
\]
Fix $0<\eps_0\le1$ with $\eps_0r_H<1/2$ and $\eps_0r_K<1/2$.
The linear perturbation has the finite expansion
\[
 L_\eps-L_0=\sum_{j=1}^{n-1}\eps^jP_j,\qquad
 P_jw=n\binom{n-1}{j}
       \frac{K_0^{n-1-j}\wedge\alpha^j\wedge\ddc w}{\Om}.
\]
Let $E_j$ be a Sobolev multiplication upper bound for
$\|P_j\|_{H^{s+2}\to H^s}$, obtained from the fixed smooth coefficient
Fourier norms, and put
$B_{n,s}=\sum_{j=1}^{n-1}\eps_0^{j-1}E_j$.
These constants are finite and fixed before solving the equation. For
$|\eps|\le\eps_0$ they give
\[
 \|L_\eps-L_0\|_{H^{s+2}\to H^s}\le|\eps|B_{n,s}.
\]
Impose also $|\eps|\le\eps_0$ and $|\eps|A_0B_{n,s}\le1/2$,
compatible with a nonzero $\eps$.
The Neumann series for $I+L_0^{-1}(L_\eps-L_0)$ gives an inverse
$L_\eps^{-1}:H^s_0\to H^{s+2}_0$ of norm at most
$A=2A_0$.

Set
\[
 Q_\eps(w)=\sum_{r=2}^n\binom nr
                   \frac{k^{n-r}\wedge(\ddc w)^r}{\Om}.
\]
On the unit ball of $H^{s+2}$, for a fixed finite $C_{n,s}$,
\begin{align}
 \|Q_\eps(w)\|_{H^s}&\le C_{n,s}\|w\|_{H^{s+2}}^2,\label{eq:quadratic}\\
 \|Q_\eps(w)-Q_\eps(z)\|_{H^s}&\le C_{n,s}
  (\|w\|_{H^{s+2}}+\|z\|_{H^{s+2}})\|w-z\|_{H^{s+2}}.
 \nonumber
\end{align}
The constant is uniform for $|\eps|\le\eps_0$ and is obtained by summing
the finite wedge coefficients, the fixed coefficient norms of
$k^{n-r}/\Om$ and the Sobolev algebra constant.
These are Sobolev algebra estimates: Fourier convolution and
Cauchy--Schwarz use
$\sum_{m\in\mathbb Z^{2n}}(1+|m|^2)^{-s}<\infty$ for $s>n$.
There is no mean obstruction in the nonlinear term. For $r\ge2$,
\[
 \int_Y k^{n-r}\wedge(\ddc w)^r
 =\int_Yw\,\ddc(k^{n-r})\wedge(\ddc w)^{r-1}=0,
\]
first for Fourier polynomials and then by Sobolev approximation.
The same Fourier estimate supplies a finite $D_{n,s}$ with
\[
 \|\ddc w\|_{L^\infty,K_0,\mathrm{op}}
                         \le D_{n,s}\|w\|_{H^{s+2}}.
\]
Choose
\begin{equation}\label{eq:radius}
 0<R\le\min\left\{1,\frac1{4AC_{n,s}},\frac1{4D_{n,s}}\right\},
 \qquad \delta=\frac{R}{2A}.
\end{equation}

For $g\in H^s_0$ with $\|g\|_{H^s}\le\delta$, the map
\[
 w\longmapsto L_\eps^{-1}(g-Q_\eps(w))
\]
maps the closed radius $R$ ball in $H^{s+2}_0$ into itself, with
image norm at most $A\delta+AC_{n,s}R^2\le3R/4$ and Lipschitz constant
at most $2AC_{n,s}R\le1/2$. Successive approximation gives a solution of
$L_\eps w+Q_\eps(w)=g$. It satisfies
\[
 k+\ddc w\ge\tfrac14K_0.
\]
The Sobolev embedding gives $w\in C^2$, so this is an equality of
ordinary continuous positive volume forms.

For the data in Theorem~\ref{thm:main}, put $F=cf$ with
$c=M/\int_Yf\Om$. Then $g=F-f_\eps$ has zero mean and satisfies the
required norm bound. Expanding the $n$th power gives
\begin{equation}\label{eq:upstairs-equation}
 (k+\ddc w)^n=(f_\eps+L_\eps w+Q_\eps(w))\Om=F\Om.
\end{equation}
The construction permits full variable perturbations. For example, in
periodic real coordinates on all elliptic factors,
\[
 g=\prod_{j=1}^n\cos(2\pi x_j)
     +\prod_{j=1}^n\cos(2\pi y_j)
\]
has zero mean. For a sufficiently small nonzero $t$, the density
$f=f_\eps+tg$ is positive, obeys \eqref{eq:data-condition}, and has
$c=1$. The prescribed equation has a nonconstant solution and treats
the full degree $n$ Monge--Amp\`ere polynomial.

\section{Bounded products and uniqueness}\label{sec:rough}

We use the following precise local facts. For decreasing locally bounded
psh functions $a_{0,j},\ldots,a_{r,j}$ with locally bounded psh limits,
both
\[
 \ddc a_{1,j}\wedge\cdots\wedge\ddc a_{r,j},\qquad
 a_{0,j}\ddc a_{1,j}\wedge\cdots\wedge\ddc a_{r,j}
\]
converge weakly to the corresponding products of their limits.
These are \cite[Theorems 2.1 and 2.4]{BT82}. Bounded recursive positivity
is also proved in \cite[Proposition 2.9]{BT76}; the locally uniform
continuous-potential result there is not used for discontinuous limits.
Their proofs in \cite{BT82} include the
top degree weighted case. We also reconstruct the weighted step below
from ordinary convergence alone. Bounded local products give no mass to
pluripolar sets: the capacity definition and
\cite[Theorem 6.9]{BT82} give this for top degree, and positivity applied
to sums of bounded psh functions with a positive quadratic gives the
mixed and lower degree versions needed here. Changing the normalization
of $\ddc$ multiplies products and capacities by fixed positive constants
and changes neither conclusion.

\begin{lemma}\label{lem:weighted}
For locally uniformly bounded decreasing psh sequences, weighted
decreasing convergence follows from ordinary mixed decreasing
convergence in one additional complex dimension.
\end{lemma}
\begin{proof}
First consider top degree on a ball in $\mathbb C^m$. Write
$\mu_j=\ddc b_{1,j}\wedge\cdots\wedge\ddc b_{m,j}$, with limit
$\mu$, and let the weight $a_j$ decrease to $a$. On a relatively
compact ball all the weights lie between constants $L$ and $U$.
Choose $0<r<R$ with $\log r^2<L-1$ and $\log R^2>U+1$.
On its product with $\{|w|<R\}$ put
\[
 v_j(z,w)=\max\{a_j(z),\log|w|^2\},\qquad
 v(z,w)=\max\{a(z),\log|w|^2\}.
\]
These are locally bounded psh functions decreasing to $v$, including
at $w=0$. Ordinary mixed decreasing convergence gives
\[
 T_j=\mu_j\wedge\ddc v_j\longrightarrow
 T=\mu\wedge\ddc v.
\]
The defining bounded-product recursion writes $T=\ddc(v\mu)$.
Since $\mu$ has top base degree, only the vertical $w$ derivatives
survive. For each fixed $z$ the resulting vertical measure is the
uniform probability on $|w|=\exp(a(z)/2)$: the radial derivative of
$\max\{t,\log|w|^2\}$ jumps from $0$ to $2/|w|$, with Euclidean
Laplacian mass $4\pi$, while our $\ddc$ is the planar Laplacian
measure divided by $4\pi$. The identical formula holds for $T_j$.
This follows by distributional testing and Fubini against the bounded
Borel function $v$; continuity of $a$ is not required.

Choose a smooth radial compactly supported test $\vartheta(w)$ in
$r<|w|<R$ equal to $\log|w|^2$ near every circle between
$\exp(L/2)$ and $\exp(U/2)$. For a smooth compactly supported base
test $\xi$ the identities are
\[
 \langle T_j,\xi\vartheta\rangle=\int\xi a_j\,\mathrm{d}\mu_j,
 \qquad \langle T,\xi\vartheta\rangle=\int\xi a\,\mathrm{d}\mu.
\]
The weak convergence proves the weighted assertion. For lower degree
currents, wedge with fixed flat psh quadratics to reach top degree;
polarization of those constant test factors recovers all current
coefficients. A quasi-psh weight is handled by adding a fixed local
psh quadratic and subtracting its smooth contribution. Uniformly
converging smooth coefficient forms may be replaced by their limits
using the uniform local product mass bounds.
\end{proof}

For a bounded canonical $k$-psh function $W$, choose locally a fixed
positive quadratic $q$ with $\ddc q\ge k$. Then $W+q$ is locally
bounded psh. Expanding relative to $\ddc(W+q)$ defines
$(k+\ddc W)^n$ by the stated products and smooth coefficient forms,
independently of $q$. This also explains its non-pluripolar property.

\begin{lemma}\label{lem:mass}
Every bounded genuine $k$-psh function $W$ satisfies
\[
 \int_Y(k+\ddc W)^n=\int_Yk^n=M.
\]
The same mass formula holds for every bounded genuine $H$-psh function.
\end{lemma}
\begin{proof}
Expand the $n$th power into terms $k^{n-r}\wedge(\ddc W)^r$.
For $r\ge1$, bounded-product recursion and compact distributional
Stokes give
\[
 \int_Yk^{n-r}\wedge(\ddc W)^r
 =\int_Y W\,\ddc(k^{n-r})\wedge(\ddc W)^{r-1}=0.
\]
This can also be obtained from the decreasing convolutions below and
unweighted product convergence; each smooth integral vanishes by the
same identity. Subtracting or adding the bounded smooth $v_0$
interchanges the $H$ and $k$ backgrounds. The identical calculation
for $H$, whose powers are also $\ddc$-closed, gives
$\int_YH^n=\int_Yk^n$.
\end{proof}

\begin{proposition}\label{prop:uniqueness}
Let $V$ be the solution of Section~\ref{sec:existence}, so
$k+\ddc V\ge\beta K_0$ with $\beta=1/4$.
If $U$ is any bounded canonical $k$-psh function with
$(k+\ddc U)^n=(k+\ddc V)^n$, then $U-V$ is constant.
\end{proposition}
\begin{proof}
On the flat torus convolve with a nonnegative radial kernel of radius
$r_j\downarrow0$. Set
\[
 U_j=U*\varrho_j,\quad V_j=V*\varrho_j,\quad
 k_j=k*\varrho_j,\quad
 A_j=k_j+\ddc U_j,\quad B_j=k_j+\ddc V_j.
\]
Translation preserves positive currents, so $A_j\ge0$ and
$B_j\ge\beta K_0$. Moreover $k_j\to k$ smoothly and
$\ddc k_j=0$, while $V_j\to V$ in $C^2$.
The convolved torsion remains pulled back from the first two factors,
so $\mathrm{d}k_j\wedge\dc k_j=0$ as well.

To check the monotone hypotheses, bound $\ddc U$ and $\ddc V$
below by a fixed negative flat form. Locally add the same positive
quadratic to each function. The spherical averages of a subharmonic
function increase with their radius, as do their radial weighted
averages. Consequently for one fixed $b\ge0$,
\[
 U_j+br_j^2\downarrow U,\qquad V_j+br_j^2\downarrow V
\]
in canonical values. The constants are the convolution correction of
the quadratic. After the local quadratic shifts, these are precisely
locally uniformly bounded decreasing psh sequences. The constants
cancel in their differences and do not change their Hessians.

Put $\varphi_j=U_j-V_j$ and
$T_j=\sum_{r=0}^{n-1}A_j^r\wedge B_j^{n-1-r}$. Smooth algebra gives
\[
 A_j^n-B_j^n=\ddc\varphi_j\wedge T_j.
\]
All $A_j,B_j$ have the same $\mathrm{d}$ and $\dc$ derivatives as
$k_j$. The Leibniz rule for $\ddc$ of each mixed monomial in $T_j$
therefore contains only factors $\ddc k_j$ or
$\mathrm{d}k_j\wedge\dc k_j$, both zero. Thus $\ddc T_j=0$.
Integrate
$\ddc(\varphi_j^2)=2\varphi_j\ddc\varphi_j+
2\mathrm{d}\varphi_j\wedge\dc\varphi_j$ against $T_j$.
Compact Stokes yields
\begin{equation}\label{eq:energy}
 -\int_Y\varphi_j(A_j^n-B_j^n)
 =\int_Y\mathrm{d}\varphi_j\wedge\dc\varphi_j\wedge T_j
 \ge\beta^{n-1}\int_Y\mathrm{d}\varphi_j\wedge\dc\varphi_j\wedge K_0^{n-1}.
\end{equation}
Only $\ddc T_j=0$ is used; $\mathrm{d}T_j$ need not vanish.

Here the left side actually tends to zero. To see each required limit,
expand $A_j$ locally relative to the positive closed current
$P_j=\ddc(U_j+br_j^2+q)$ and the smooth forms
$k_j-\ddc q$. Weighted monotone convergence applies to
$U_jA_j^n$ by adding and then subtracting $q+br_j^2$.
Unweighted convergence and uniform convergence of $V_j$ apply to
$V_jA_j^n$. The terms $U_jB_j^n$ converge by dominated convergence:
$U_j$ are uniformly bounded and converge pointwise, and $B_j^n$
have uniformly converging continuous densities. The terms $V_jB_j^n$
converge by $C^2$ convergence. A finite partition of unity patches
these local product limits. Thus the limit of the left side is
\[
 -\int_Y(U-V)\big((k+\ddc U)^n-(k+\ddc V)^n\big)=0.
\]
The flat Dirichlet energy on the right of \eqref{eq:energy} therefore
tends to zero. Bounded pointwise convergence gives
$\varphi_j\to U-V$ in $L^2$. Fourier Poincar\'e then makes
$U-V$ constant almost everywhere, and the canonical quasi-psh
representatives give this identity at every point.
\end{proof}

Lemma~\ref{lem:mass} first forces equality of the constants in any
two equations with the same nonzero datum. Proposition~\ref{prop:uniqueness}
then compares a rough solution with the constructed positive solution;
the competitor is not required to lie in the contraction ball or to be smooth.

\section{Disk descent, measure factors and continuity}

Every holomorphic disk into the cusp lifts holomorphically. If its
components satisfy $y^2=x^3$ and $x\not\equiv0$, define $z=y/x$ off
the zeros of $x$. Then $z^2=x$, $z^3=y$, and $|z|^2=|x|$;
the bounded removable singularities extend $z$ across those zeros.
If $x\equiv0$, also $y\equiv0$, and the lift is constant.
All other elliptic components are unchanged. This proves disk lifting
for every disk into $X$.

Since $\pi$ is a homeomorphism, the function
\[
 u=(v_0+w)\circ\pi^{-1}-\sup_Y(v_0+w)
\]
is continuous. For any disk $f$ and its holomorphic lift $\widetilde f$,
\[
 \ddc(u\circ f+p_f)
 =\widetilde f^*(H+\ddc(v_0+w))\ge0.
\]
It is therefore genuinely $h$-psh, with finite canonical values, and
\eqref{eq:upstairs-equation} gives the desired pushed-forward equation.
Conversely a bounded genuine $h$-psh function pulls back to a bounded
genuine $H$-psh function by disk composition. Pushforward by the Borel
homeomorphism $\pi$ is injective on finite measures: for Borel $B\subset Y$,
$\mu(B)=(\pi_*\mu)(\pi(B))$. Thus any bounded downstairs equation
with datum $\pi_*(f\Om)$ lifts to exactly
$(H+\ddc(u\circ\pi))^n=cf\Om$, with the same constant.
Subtracting $v_0$ reduces
its equation to the positive $k$ background of Section~\ref{sec:rough}.
Thus the upstairs uniqueness covers the entire downstairs ordinary type.

The datum has no mass on the singular locus, whose inverse image $D$
has smooth $\Om$-measure zero. A local nontrivial psh pole witness pulls
back through this finite bijective normalization to a nontrivial psh
witness; such witnesses are locally integrable and their pole sets are
smooth-volume-null \cite[Chapter I, (5.3) and Theorem (4.17)]{Dem}.
Consequently $\nu$ is non-pluripolar.

The generic degree of $\pi$ is one. On $Y\setminus D$ set
$J=H^n/\Om>0$. Under the biholomorphic identification with $X_{\rm reg}$,
the density is exactly $f/J$ relative to $h^n$.
Near the cusp, where $\alpha=0$,
\[
 \pi^*h^n=n!\rho\kappa_1\wedge\cdots\wedge\kappa_n,
 \qquad \Om=n!\lambda\kappa_1\wedge\cdots\wedge\kappa_n.
\]
Thus on this punctured cusp chart the density transformation is
\[
 \nu=\frac{f\lambda}{\rho}\,h^n.
\]
Here $\rho=4|z|^2+9|z|^4$. The reciprocal factor has order
$|z|^{-2}$; the measure is the pushforward of a finite volume with
continuous density and carries no exceptional mass. The equation constant is forced
by $M=c\nu(X)$, with no missing degree or factorial.

For the continuity estimate, let $d=z-z'$ and $a'=z+z'$. Then
\[
 z^2-z'^2=da',\qquad
 z^3-z'^3=\frac{d(3a'^2+d^2)}4.
\]
If $|a'|\ge|d|/4$, the first difference has magnitude at least
$|d|^2/4$, which controls $|d|^3$ on a bounded chart. Otherwise the
second has magnitude at least $13|d|^3/64$. Hence
\[
 |z-z'|\le C\bigl(|z^2-z'^2|+|z^3-z'^3|\bigr)^{1/3}.
\]
Combining this with the ordinary other factor coordinates and the
$C^1$ regularity of $v_0+w$ proves the asserted H\"older bound for the
full potential across the cusp.

\section{All refinements and the scope}\label{sec:refinements}

Every smooth resolution $Z\to X$ factors through $Y$. Locally, pull back
the cusp coordinates $x,y$ and extend $z=y/x$ from the dense regular
preimage. The estimate $|z|^2=|x|$ and normality of the smooth source
give its holomorphic extension; the resulting lifts glue uniquely.
The factor $\mu:Z\to Y$ is a proper birational modification.

For every bounded canonical $H$-psh $W$, the pullback $W\circ\mu$
has bounded local psh shifts. Its bounded products give no mass to the
analytic exceptional locus. Outside that locus, $\mu$ is biholomorphic
and the products agree by ordinary local change of variables. Therefore
\[
 \mu_*\MA_{\mu^*H}(W\circ\mu)=\MA_H(W).
\]
The same reasoning applies to every common refinement. None of these
pullback backgrounds is declared strictly positive. The constructed
solution on each model is literally a composition with $\mu$, so it
is constant on every compact $\mu$-fiber with its finite canonical value.
This proves the measure independence and the actual descent assertion.

The specified ordinary type is preserved and recovered under these
surjective pullbacks. For functions $a,b$ on $X$, the inequality
$a\le b+C$ holds downstairs if and only if
$a\circ\sigma\le b\circ\sigma+C$ holds on a resolution: the reverse
implication follows by choosing a preimage of each point. Apply this
in both directions. In particular boundedness relative to the zero
reference, and hence exactly the stated bounded ordinary type, is
independent of the chosen resolution or common refinement.

The examples exist in every complex dimension $n\ge2$ with a nonnormal cusp
divisor and a genuinely nonclosed background. The Fourier primitive on
$E_1$ repairs the reference; the prescribed density equation is a nonlinear
equation on all $2n$ real variables. The result concerns bounded ordinary
type and an open positive Sobolev density family on this geometry. It makes
no claim for arbitrary prescribed pole types, arbitrary Borel data, arbitrary
unibranch spaces, or arbitrary nonclosed torsion. Those parts of P23 remain open.

These spaces also differ globally from projective cones. Every holomorphic
map $\mathbb P^1\to X$ lifts locally by the disk calculation, and the lifts
glue uniquely to a map into $Y$. Each flat holomorphic one-form
$\mathrm{d}z_j$ pulls back to zero on $\mathbb P^1$. Thus the lifted map
has zero derivative and is constant. The space $X$ contains no nonconstant
rational curves, whereas a nontrivial projective cone contains its ruling
projective lines.

\begin{thebibliography}{9}
\bibitem{BT76} E. Bedford and B. A. Taylor,
\emph{The Dirichlet problem for a complex Monge--Amp\`ere equation},
Invent. Math. \textbf{37} (1976), 1--44.
\href{https://doi.org/10.1007/BF01418826}{doi:10.1007/BF01418826}.
\bibitem{BT82} E. Bedford and B. A. Taylor,
\emph{A new capacity for plurisubharmonic functions},
Acta Math. \textbf{149} (1982), 1--40.
\href{https://doi.org/10.1007/BF02392348}{doi:10.1007/BF02392348}.
\bibitem{Dem} J.-P. Demailly,
\emph{Complex Analytic and Differential Geometry}, version of 21 June 2012.
\href{https://www-fourier.univ-grenoble-alpes.fr/~demailly/manuscripts/agbook.pdf}{Author's manuscript}.
\end{thebibliography}
\end{document}
